
\documentclass[10pt]{article}
\usepackage[utf8]{inputenc}
\usepackage[T1]{fontenc}
\usepackage{amsmath, amssymb, xcolor, geometry, enumitem, multicol, tcolorbox}

\definecolor{mainblue}{RGB}{0, 102, 204}
\geometry{a4paper, margin=1.2cm, top=1cm, bottom=3cm, footskip=1.2cm}

\begin{document}
\enlargethispage{2.5cm}

% Modern Header
\begin{tcolorbox}[colback=mainblue!5, colframe=mainblue, sharp corners, boxrule=0.5pt]
    \small \textbf{Unit:} Unit 2: Exponents & Radicals \hfill \textbf{Name:} \rule{3cm}{0.4pt} \\
    \small \textbf{Topic:} Exponents \hfill \textbf{Date:} \rule{2cm}{0.4pt} \\
    \centering \textbf{\large Intro (Level 0)}
\end{tcolorbox}

\vspace{0.2cm}

% MCQ Section
\begin{multicols}{2}
\begin{enumerate}[label=\textbf{Q\arabic*.}, leftmargin=*, itemsep=6pt, parsep=0pt]

  \item \small What is the value of 2^3?
  \begin{enumerate}[label=(\Alph*), leftmargin=0.6cm, nosep, itemsep=2pt]
    \item 6
    \item 8
    \item 9
    \item 10
    \item 12
  \end{enumerate}

  \item \small What is the result of 3^2?
  \begin{enumerate}[label=(\Alph*), leftmargin=0.6cm, nosep, itemsep=2pt]
    \item 5
    \item 6
    \item 8
    \item 9
    \item 10
  \end{enumerate}

  \item \small Simplify the expression: 4^1
  \begin{enumerate}[label=(\Alph*), leftmargin=0.6cm, nosep, itemsep=2pt]
    \item 2
    \item 4
    \item 8
    \item 16
    \item 32
  \end{enumerate}

  \item \small Evaluate the expression: 2^4
  \begin{enumerate}[label=(\Alph*), leftmargin=0.6cm, nosep, itemsep=2pt]
    \item 8
    \item 10
    \item 12
    \item 14
    \item 16
  \end{enumerate}

  \item \small Simplify the expression: 5^0
  \begin{enumerate}[label=(\Alph*), leftmargin=0.6cm, nosep, itemsep=2pt]
    \item 0
    \item 1
    \item 5
    \item 10
    \item 25
  \end{enumerate}

  \item \small What is the value of (2^2)^3?
  \begin{enumerate}[label=(\Alph*), leftmargin=0.6cm, nosep, itemsep=2pt]
    \item 16
    \item 24
    \item 32
    \item 64
    \item 128
  \end{enumerate}

  \item \small Simplify the expression: (3^2)^2
  \begin{enumerate}[label=(\Alph*), leftmargin=0.6cm, nosep, itemsep=2pt]
    \item 9
    \item 16
    \item 25
    \item 27
    \item 81
  \end{enumerate}

  \item \small What is the result of 2^0?
  \begin{enumerate}[label=(\Alph*), leftmargin=0.6cm, nosep, itemsep=2pt]
    \item 0
    \item 1
    \item 2
    \item 4
    \item 8
  \end{enumerate}

\end{enumerate}
\end{multicols}

\vspace{-0.2cm}
\hrule
\vspace{0.2cm}
\textbf{\small Open Ended Questions (FRQ)}
\begin{enumerate}[label=\textbf{Q\arabic*.}, start=9, itemsep=5pt, topsep=0pt]

  \item \small Simplify the expression: 2^3 * 2^4 \vspace{1.2000000000000002cm}

  \item \small Simplify the expression: (2^2)^3 * 2^2 \vspace{1.2000000000000002cm}

\end{enumerate}

\vfill

% Chef's Tips Footer
\begin{tcolorbox}[colback=blue!5, colframe=mainblue!50, title=\small \textbf{Chef's Tips}, fonttitle=\bfseries, bottom=1mm]
\begin{itemize}[noitemsep, topsep=2pt, leftmargin=0.5cm]
  \item \scriptsize Tip 1: Remember that any number raised to the power of 0 is 1\item \scriptsize Tip 2: To multiply two powers with the same base, add the exponents\item \scriptsize Tip 3: To divide two powers with the same base, subtract the exponents
\end{itemize}
\end{tcolorbox}

\end{document}
  